% ============================================================
% v2/sections_v2/conclusion.tex  おわりに（査読の指摘を受けて改訂 2026-10-27 17:00）
% ============================================================
\section{おわりに}\label{sec:conclusion}
本研究は，重ね合わせ機構\cite{noguchi2022superposition}の上で，自分の価値の経験から学んだ予測傾向がどのような条件で他者側の予測に転用されるかを調べた．
価値の信号に従って動いた経験からは，報酬も目標の入力も用いない予測学習によって価値に従う行動の見込みが生まれ，それは共有された重みを通して他者側に与えた快・不快にも反映された．
同じ軌跡で価値の信号と行動の対応だけを壊すとこの転用は大きく弱まり，他者の状態を追う力だけでは生じず，快・不快の入口の重みを写すだけでは戻らなかった．
一方，転用は他者の行動の予測を一般に正確にはせず，他者固有の好みの自発的な表現や推定も保証しなかった．
自分の経験に基づく予測は他者予測の出発点になりうるが，それだけでは他者の理解に足りない．
これらの条件の切り分けは，自己位置の追跡の前提を満たさなかった経験80\%条件のモデルで得たものであり，前提を満たすモデルでの確認が残る．
